\documentclass[11pt]{article}

\usepackage{bm}
\usepackage[utf8]{inputenc}
\usepackage[italian]{babel}
\usepackage{amsmath, amssymb}
\usepackage{amsthm}
\usepackage{graphicx}
\usepackage{booktabs}
\usepackage{geometry}
\geometry{a4paper, top=2.5cm, bottom=2.5cm, left=2.5cm, right=2.5cm}
\usepackage{caption}
\captionsetup{justification=centering}
\usepackage{float}
\usepackage[hidelinks]{hyperref}
\usepackage{subcaption}
\usepackage{verbatim}
\usepackage{enumitem}
\usepackage{yhmath}
\usepackage{derivative}
\usepackage{physics}

\theoremstyle{plain}
\newtheorem{theorem}{Teorema}
\begin{document}

\title{Forza gravitazionale nulla in un guscio sferico omogeneo}
\author{Liam Ferretti}
\maketitle

Sia $p$ un generico punto interno ad un guscio sferico omogeneo, allora la forza gravitazionale risultante è nulla.

\begin{proof}
    Sia $\vec{q}$ un punto appartenente alla superficie di un guscio sferico omogeneo sottile, sia $\vec{p}$ un punto interno al guscio, allora è possibile senza perdere di generalità orientare una terna destrorsa di assi cartesiani in modo tale che il punto $\vec{p}$ giaccia sul semi asse $z$ positivo, segue che $\vec{p}$ e $\vec{q}$ sono definiti come segue:
    \begin{equation}
        \vec{p} := (0, 0, z_0), \quad \quad \vec{q} := (R\sin\theta\cos\varphi, R\sin\theta\sin\varphi, R\cos\theta)
    \end{equation}
    nella quale:
    \begin{equation}
        0 \leq z_0 < R, \quad \quad \theta \in [0, \pi], \quad \quad \varphi \in [0, 2\pi]
    \end{equation}
    Allora è possibile definire z come la distanza tra p e q:
    \begin{equation}
        \vec{z} = \vec{p} - \vec{q} = (- R\sin\theta\cos\varphi, - R\sin\theta\sin\varphi, z_0 - R\cos\theta)
    \end{equation}
    \begin{equation}
        \implies |z| = \sqrt{R^2 + z_0^2 - 2Rz_0\cos\theta} 
    \end{equation}
    Sapendo che:
    \begin{equation}
        dA = R^2\sin\theta d\theta d\varphi \implies dM = \sigma dA = \sigma R^2 \sin\theta d\theta d\varphi
    \end{equation}
    Allora vale che:
    \begin{equation}
        dF = G\frac{mdM}{|z|^3}\vec{z}
    \end{equation}
    allora è possibile studiare le componenti $x, y, z$ nel seguente modo:
    \begin{equation}
        \begin{cases}
            \displaystyle dF_x = \frac{-Gm\sigma R^3\sin^2\theta\cos\varphi}{(R^2 + z_0^2 - 2Rz_0\cos\theta)^{3/2}}d\theta d\varphi \\[1.2em]
            \displaystyle dF_y = \frac{-Gm\sigma R^3\sin^2\theta\sin\varphi}{(R^2 + z_0^2 - 2Rz_0\cos\theta)^{3/2}}d\theta d\varphi \\[1.2em]
            \displaystyle dF_z = \frac{Gm\sigma R^2 (z_0 - R\cos\theta)\sin\theta}{(R^2 + z_0^2 - 2Rz_0\cos\theta)^{3/2}}d\theta d\varphi
        \end{cases}
    \end{equation}
    Allora integrando la componente x della forza, segue che:
    \begin{align}
        F_x & = \int dF_x \\
            & = \int_{0}^{\pi}\int_{0}^{2\pi} \frac{-Gm\sigma R^3\sin^2\theta\cos\varphi}{(R^2 + z_0^2 - 2Rz_0\cos\theta)^{3/2}}d\theta d\varphi \\
            & = -Gm\sigma R^3 \int_{0}^{\pi} \frac{\sin^2\theta}{(R^2 + z_0^2 - 2Rz_0\cos\theta)^{3/2}}d\theta \int_{0}^{2\pi} \cos\varphi d\varphi \\
            & = -Gm\sigma R^3 \int_{0}^{\pi} \frac{\sin^2\theta}{(R^2 + z_0^2 - 2Rz_0\cos\theta)^{3/2}}d\theta \cdot \left[\eval{\sin\varphi}_{0}^{2\pi}\right] \\
            & = -Gm\sigma R^3 \int_{0}^{\pi} \frac{\sin^2\theta}{(R^2 + z_0^2 - 2Rz_0\cos\theta)^{3/2}}d\theta \cdot 0 \\
            & = 0
    \end{align}
    Analogamente integrando la componente lungo y segue che:
    \begin{align}
        F_y & = \int dF_y \\
            & = \int_{0}^{\pi}\int_{0}^{2\pi} \frac{-Gm\sigma R^3\sin^2\theta\sin\varphi}{(R^2 + z_0^2 - 2Rz_0\cos\theta)^{3/2}}d\theta d\varphi \\
            & = -Gm\sigma R^3 \int_{0}^{\pi} \frac{\sin^2\theta}{(R^2 + z_0^2 - 2Rz_0\cos\theta)^{3/2}}d\theta \int_{0}^{2\pi} \sin\varphi d\varphi \\
            & = -Gm\sigma R^3 \int_{0}^{\pi} \frac{\sin^2\theta}{(R^2 + z_0^2 - 2Rz_0\cos\theta)^{3/2}}d\theta \cdot \left[\eval{-\cos\varphi}_{0}^{2\pi}\right] \\
            & = -Gm\sigma R^3 \int_{0}^{\pi} \frac{\sin^2\theta}{(R^2 + z_0^2 - 2Rz_0\cos\theta)^{3/2}}d\theta \cdot 0 \\
            & = 0
    \end{align}

    Integrando invece la componente $z$, si ottiene:
    \begin{align}
        F_z & = \int dF_z \\
            & = \int_{0}^{\pi} \int_{0}^{2\pi} \frac{Gm\sigma R^2 (z_0 - R\cos\theta)\sin\theta}{(R^2 + z_0^2 - 2Rz_0\cos\theta)^{3/2}}d\theta d\varphi \\
            & = 2\pi Gm\sigma R^2 \cdot \int_{0}^{\pi} \frac{(z_0 - R\cos\theta)\sin\theta}{(R^2 + z_0^2 - 2Rz_0\cos\theta)^{3/2}}d\theta
    \end{align}

    Imponendo la prima sostituzione:
    \begin{equation}
        u = \cos\theta \implies du = -\sin\theta d\theta
    \end{equation}
    I vecchi estremi di integrazione cambiano come segue:
    \begin{align}
        \theta = 0 & \implies u = \cos(0) = 1 \\
        \theta = \pi & \implies u = \cos(\pi) = -1 
    \end{align}

    Sostituendo nell'integrale si ha:
    \begin{align}
        F_z & = - 2\pi Gm\sigma R^2 \cdot \int_{1}^{-1} \frac{z_0 - Ru}{(R^2 + z_0^2 - 2Rz_0 u)^{3/2}}du \\
            & = 2\pi Gm\sigma R^2 \cdot \int_{-1}^{1} \frac{z_0 - Ru}{(R^2 + z_0^2 - 2Rz_0 u)^{3/2}}du 
    \end{align}

    Per calcolare esplicitamente la primitiva dell'integrale indefinito:
    \begin{equation}
        J = \int \frac{z_0 - Ru}{(R^2 + z_0^2 - 2Rz_0 u)^{3/2}}du
    \end{equation}
    Imponendo t come segue:
    \begin{equation}
        t = R^2 + z_0^2 - 2Rz_0 u \implies dt = -2Rz_0 \, du \implies du = -\frac{dt}{2Rz_0}
    \end{equation}
    segue che:
    \begin{equation}
        2Rz_0 u = R^2 + z_0^2 - t \implies Ru = \frac{R^2 + z_0^2 - t}{2z_0}
    \end{equation}
    Il numeratore diventa quindi:
    \begin{equation}
        z_0 - Ru = z_0 - \frac{R^2 + z_0^2 - t}{2z_0} = \frac{2z_0^2 - R^2 - z_0^2 + t}{2z_0} = \frac{z_0^2 - R^2 + t}{2z_0}
    \end{equation}

    Riscriviamo l'integrale $J$ in funzione di $t$ separando i termini:
    \begin{align}
        J & = \int \frac{\frac{z_0^2 - R^2 + t}{2z_0}}{t^{3/2}} \left( -\frac{dt}{2Rz_0} \right) \\
        & = -\frac{1}{4Rz_0^2} \int \left( (z_0^2 - R^2)t^{-3/2} + t^{-1/2} \right) dt
    \end{align}

    Integrando singolarmente le potenze di $t$:
    \begin{align}
        J & = -\frac{1}{4Rz_0^2} \left[ (z_0^2 - R^2) \left( \frac{t^{-1/2}}{-1/2} \right) + \frac{t^{1/2}}{1/2} \right] \\
        & = -\frac{1}{4Rz_0^2} \left[ -\frac{2(z_0^2 - R^2)}{\sqrt{t}} + 2\sqrt{t} \right] \\
        & = \frac{z_0^2 - R^2}{2Rz_0^2\sqrt{t}} - \frac{\sqrt{t}}{2Rz_0^2}
    \end{align}

    Riconducendo tutto ad un unico denominatore comune $2Rz_0^2\sqrt{t}$:
    \begin{equation}
        J = \frac{(z_0^2 - R^2) - t}{2Rz_0^2\sqrt{t}}
    \end{equation}

    Ripristiniamo la variabile $u$ sostituendo nuovamente $t = R^2 + z_0^2 - 2Rz_0 u$ al numeratore:
    \begin{align}
        (z_0^2 - R^2) - (R^2 + z_0^2 - 2Rz_0 u) & = z_0^2 - R^2 - R^2 - z_0^2 + 2Rz_0 u \\
                                            & = 2Rz_0 u - 2R^2 = 2R(z_0 u - R)
    \end{align}

    da cui segue che:
    \begin{equation}
        J = \frac{2R(z_0 u - R)}{2Rz_0^2 \sqrt{R^2 + z_0^2 - 2Rz_0 u}} = \frac{z_0 u - R}{z_0^2 \sqrt{R^2 + z_0^2 - 2Rz_0 u}}
    \end{equation}

    Tornando all'integrale definito $F_z$:
    \begin{equation}
        F_z = 2\pi Gm\sigma R^2 \cdot \left[ \frac{z_0 u - R}{z_0^2 \sqrt{R^2 + z_0^2 - 2Rz_0 u}} \right]_{-1}^{1}
    \end{equation}

    Raccogliendo la costante $z_0^2$ a denominatore e sostituendo i limiti superiore e inferiore:
    \begin{align}
        F_z & = \frac{2\pi Gm\sigma R^2}{z_0^2} \cdot \left[ \frac{z_0(1) - R}{\sqrt{R^2 + z_0^2 - 2Rz_0(1)}} - \frac{z_0(-1) - R}{\sqrt{R^2 + z_0^2 - 2Rz_0(-1)}} \right] \\
         & = \frac{2\pi Gm\sigma R^2}{z_0^2} \cdot \left[ \frac{z_0 - R}{\sqrt{(z_0 - R)^2}} - \frac{-z_0 - R}{\sqrt{(z_0 + R)^2}} \right]
    \end{align}
    \begin{align}
        \sqrt{(z_0 - R)^2} = |z_0 - R| & = R - z_0 \\
        \sqrt{(z_0 + R)^2} = |z_0 + R| & = z_0 + R
    \end{align}

    \begin{align}
        F_z & = \frac{2\pi Gm\sigma R^2}{z_0^2} \cdot \left[ \frac{z_0 - R}{R - z_0} - \frac{-(z_0 + R)}{z_0 + R} \right] \\
            & = \frac{2\pi Gm\sigma R^2}{z_0^2} \cdot \left[ -1 - (-1) \right] \\
            & = \frac{2\pi Gm\sigma R^2}{z_0^2} \cdot 0 \\
            & = 0
    \end{align}

    Da cui segue che:
    \begin{equation}
        F = (0, 0, 0)
    \end{equation}
    Nel caso in cui $z_0 = 0$ lo sviluppo per $F_z$ non è più valido poichè si dividerebbe per 0, ma tale caso si risolve banalemente per simmetria.

    In conclusione, la forza totale agente sul punto interno al guscio sferico è nulla.
\end{proof}

\end{document}